\section{Literature Review}  \label{sec:lit_review}
% {\bf Related Literature.} 

{\bf Roads and House Prices.} \cite{gonzalez2016paving} studies the effects of paving roads in Acayucan, Mexico. The government identified 56 neighborhoods that needed paved roads. Of these, 28 were randomly selected for paving. About 1,000 people were surveyed in these neighborhoods before paving (2006) and after (2009), although 11 of the 28 treatment groups were still in the process of being paved at the time of the post-treatment survey.  It finds a 17\% increase in property values as measured by professional appraisals, a 28\% increase in homeowner-estimated property values, a 36\% increase in rents, and a 72\% increase in vacant land values, along with effects on a few non-housing outcomes.  Our study evaluates cuts in road maintenance rather than a switch from unpaved to paved roads, and our Ohio geography contrasts with that of a developing nation.  Our housing values are not based on professional appraisals or homeowner valuation but on actual sales transactions; and the fact that we observe sales transactions every year lets us analyze pre- and post-treatment trends in house prices over a long timeframe rather than a one-time change in house values. 

% Another study, \cite{wong2017local}, studies the effect of local political reform in rural chinese villages on the quality of road projects and finds, using fixed 

Recent empirical work by \cite{gertler2024road} studies highway maintenance in Indonesia and finds that better road quality energizes local labor markets and increases housing prices. Although both our study and \cite{gertler2024road} examine road maintenance and its effects on housing prices, our focus on local roads in a developed country contrasts with their focus on national highways in a developing country. We also differ in our road quality measurement as they use International Roughness Index (IRI) data collected by the Indonesian Ministry of Public Works, while we develop our own road quality measure using high-resolution satellite imagery processed through a fine-tuned vision model. Most importantly, our identification strategies differ fundamentally. We exploit close elections to renew local road maintenance levies in a sharp regression discontinuity design, while \cite{gertler2024road} employs an instrumental variables approach based on Indonesia's two-stage budgeting process. Their strategy instruments for district-level road quality using provincial-level road maintenance spending, arguing that central government budget allocations to provinces are exogenous to district-level economic shocks. However, their approach faces challenges. First, they proxy for the exogenous and unavailable financial budget data using total kilometers of roads upgraded at the provincial level, which may itself be endogenous: provinces experiencing economic shocks could see changes in both upgrade intensity and economic outcomes, violating the exclusion restriction. Second, upgrade costs per kilometer may vary substantially across terrain and location, making kilometers upgraded a noisy proxy for financial resources. In contrast, our sharp RD design avoids these concerns by exploiting the quasi-random assignment of treatment around the election threshold and focusing on close elections where the margin of victory is uncorrelated with potential outcomes, providing a cleaner path to identification.


% We study the maintenance of local roads in a developed nation, whereas \cite{gertler2024road} focuses on national highways in a developing nation.  Our identification strategy also differs, as we exploit close elections to renew road maintenance levies, while \cite{gertler2024road} uses an instrumental variables approach based on the timing of routine maintenance cycles.  

Theoretical work by \cite{rioja2003} develops a dynamic general equilibrium model to study the optimal amount of road maintenance, which is found to depend on the size of new infrastructure investments and the productivity of infrastructure. \cite{rioja2003} solves and parameterizes the model, finding that increasing funding away from new infrastructure and toward the maintenance of existing infrastructure decreases the depreciation rate of existing infrastructure, increases the stock of existing infrastructure, and increases economic output and consumption levels. 
   
% Another study, \cite{chaurey2022}, assesses the impact of infrastructure maintenance in India. As their research notes, the Indian national government identified 17 major states as the poorest and made an index of ``backwardness'' for districts in these states. The 115 most backward districts were awarded 450 million rupees and were given discretion on how to spend these funds to ``improve and maintain'' or ``make complementary investments'' to existing infrastructure.  Money could be used to widen roads, add lanes, add electricity transmission and distribution infrastructure, build new road links to markets, and build bridges, for example.  The treatment is therefore a mixture of road and non-road spending and of maintenance and new construction, especially, presumably, in districts that did not already have paved roads and electricity, where \cite{chaurey2022} finds the largest effects. 
% There were two other village-level infrastructure programs in place at the time, but \cite{chaurey2022} argues that the estimates are unaffected by these programs because treatment was based on population, not backwardness, and the inclusion of population as a covariate ensures that estimates are unaffected. 
% We worry that some districts could satisfy both the population and ``backwardness'' criteria and that the inclusion of covariates in regression discontinuity does not ``control'' for anything but merely serves to increase the precision of treatment effect estimates \citep{lee_lemieux2010}. 
% \cite{chaurey2022} does not study the effects of infrastructure spending on real estate values like we do, instead providing causal estimates of infrastructure investment mainly on outcomes related to employment. 
% Unlike \cite{chaurey2022}, we focus on the effects of infrastructure spending on real estate values rather than employment outcomes.

The only other study we find on maintenance and house prices is the classic capitalization study by \cite{edel1974taxes}. It studies the effect of local public spending on median house values in the Boston area using decennial Census years from 1930 to 1970. Controlling for latitude and longitude, the tax rate, population density, tenure status, and school expenditures, its ordinary least squares regressions do not find a statistically significant link between road maintenance spending and house prices.

% {\bf Identification.} A study with an identification strategy similar to ours is \cite{cellini2010value}, which studies the effect of new capital projects for schools, funded via new local bond issues and raised by referendums. Both our study and \cite{cellini2010value} rely on broadly similar identification strategies and employ a dynamic regression discontinuity design and analyze changes in regional property values. Moreover, both papers have similar assignment mechanisms: in each case, votes in favor of or against a local referendum serves as the source of exogenous variation. Despite these similarities, a few key distinctions stand out. First, \cite{cellini2010value} looks at the impact of bond elections for additional funding on new capital infrastructure projects, while we focus on the impact of cutting renewal referendums for existing road maintenance. Second, \cite{cellini2010value} implements a fuzzy regression discontinuity approach, whereas ours is sharp. Lastly, we focus on the Intent-to-Treat (ITT) estimator, which analyzes the effect of an election on outcomes without controlling for the results of subsequent elections\footnote{Recent papers such as \cite{hsu2024dynamic} and \cite{biasi_what_2025} have focused on improving upon methods from \cite{cellini2010value}. However, these papers focus on the effects of new capital projects rather than maintaining existing capital infrastructure. Given the exogeneity of the timing of renewal elections and our focus on existing maintenance, the ITT estimator is more appropriate for our setting.} (more details in Section \ref{sec:empirical}). As \cite{cellini2010value} shows, when elections are independent, the ITT estimator is equal to the Treatment on the Treated (TOT) estimator. 

% Another paper, \cite{asher2020} studies the impact of new roads on villages in India. Like our study, its identification strategy is regression discontinuity, although ours is sharp rather than the fuzzy form.  It argues the main obstacle to identification in prior studies is that the placement of new roads is usually correlated with economic (or political) characteristics rather than exogenous. Its findings suggest this is a serious problem with the literature because, unlike prior studies, it finds no strong link between economic growth and new road placement, suggesting that the estimates of previous studies that find a link are driven by road placement in villages that are already growing. \cite{asher2020} touts its use of village-level rather than regional-level data.  We, too, look at economic outcomes at the level of village, city, and township, the most local levels of government. A surprising finding of \cite{asher2020} is that investment in transportation infrastructure does not affect village incomes, assets, or agricultural output. Its measure of assets is a village-level average of a series of binary indicators of ownership of a variety of assets, along with separate regressions for the presence of a ‘solid house’, refrigerator, and phone; whereas we study the effect of local road tax cuts on housing sale prices.  Of course, our use of a developed geography contrasts with rural villages in India.  Our efforts to achieve identification focus on the maintenance of existing roads, which avoids the endogeneity of the placement of new roads. We also find similarities with \cite{boudot2024paving}, which employs a similar identification strategy but studies new roads instead of the maintenance of existing roads and finds that villagers reward the political party responsible for spearheading the road expansion.


{\bf Road Quality.} Measuring road quality remains empirically challenging. \cite{currier2023} exploits vertical‐acceleration signals from millions of Uber trips to construct a Road Roughness Index (RRI) and shows that it moves predictably with resurfacing events. Although RRI is highly correlated with engineering benchmarks such as the International Roughness Index (IRI) and the Pavement Condition Index (PCI), each of these legacy measures has drawbacks that limit their use for large-scale economic analysis: IRI data are collected only where a costly profilometer van has recently driven, and PCI ratings depend on labour-intensive, inspector-specific visual surveys that are updated irregularly. Moreover, the Uber-based approach cannot cover streets outside the platform’s service footprint—precisely where many suburban and rural housing markets are located. In contrast, high-resolution satellite imagery is available nationwide and available at low marginal cost. By fine-tuning a vision model on a labelled road-surface dataset, we generate a consistent, wall-to-wall measure of road quality that spans urban, suburban and rural networks alike, and captures the visual cues that prospective home-buyers directly observe. 

We find another study, \cite{wong2017local}, which examines the effect of local political reform in rural Chinese villages on local infrastructure projects and finds that these reforms result in younger village leaders and higher quality roads. It measures road quality by creating a village-level road quality rating, which is based on the observations of field surveyors that use their ``road quality evaluation scheme''. They emphasize minimization of subjectivity in the evaluation process, as the surveyors are trained to use a standardized set of criteria to assess road quality, such as number of bends per 100 meters and road width. Unlike \cite{wong2017local}, which relies on surveyor-based evaluations, our road quality measure is derived from high-resolution satellite imagery processed through our fine-tuned vision model, ensuring consistency and scalability across diverse geographic regions.

% We highlight a substantial literature studying the effect of transportation infrastructure on house prices. \cite{hoogendoorn2019house} studies the effect of the opening of a tunnel on house prices in the Netherlands, noting that prior research studying transportation infrastructure in developed geographies suffers from reverse causality.  It argues that the opening of the Westerscheldtunnel is a fairly exogenous event, with natural borders that prevent contamination of results by the surrounding environment. \cite{hoogendoorn2019house} finds half the capitalized value of the tunnel occurs more than a year before the tunnel opens and argues that the exogeneity of the opening of the tunnel, along with hedonic controls, time trends and postcode fixed effects, identifies its estimates. Our data, too, is for a developed nation. One novelty of our study is how ordinary the events are that we study. While the opening of a new tunnel is important, it is rare.  Votes to renew infrastructure spending are common events in many local governments in the United States, and the quantity of road maintenance spending is regularly chosen by governments around the world, if not by voting then directly by bureaucrats. It is therefore also important to study the effects of road maintenance spending on house prices. 

% {\bf Transportation and House Prices.} We highlight a substantial literature studying the effect of transportation infrastructure on house prices. \cite{hoogendoorn2019house} studies the effect of the opening of a tunnel on house prices in the Netherlands, noting that prior research on transportation infrastructure in developed regions often suffers from reverse causality. It argues that the opening of the Westerscheldtunnel is a fairly exogenous event, with natural borders that prevent contamination of results by the surrounding environment. \cite{hoogendoorn2019house} finds that half the capitalized value of the tunnel occurs more than a year before the tunnel opens and argues that the exogeneity of the tunnel's opening, along with hedonic controls, time trends, and postcode fixed effects, identifies its estimates. Our data also pertains to a developed nation. One novelty of our study is how ordinary the events are that we study. While the opening of a new tunnel is significant, it is rare. Votes to renew infrastructure spending are common events in many local governments in the United States, and the amount of road maintenance spending is regularly determined by governments worldwide, either through voting or directly by bureaucrats. It is therefore important to study the effects of road maintenance spending on house prices.

% \cite{li2016wheels} studies the overall effect on apartment prices of new subway lines in Beijing, but the estimates may represent the net effect of competing factors. \cite{gibbons2005valuing}, studying the construction of new rail stations for the London underground and light rail services, notes that the effect on house prices captures the net effect of better access, increased crime, and increased noise pollution. \cite{levkovich2016effects} looks at the effect of highway development on house prices in the Netherlands. It separates out accessibility effects from noise pollution and increased traffic effects by looking at different neighborhoods near the highway development. Its repeat sales difference in differences model finds increased house prices from anticipation effects \citep{kohlhase1991impact}. \cite{beenstock2016hedonic} also finds anticipation effects for house prices for the development of a highway across Israel. 

{\bf Contribution.} Our paper contributes to the literature on three main fronts. 
First, while much of the literature focuses on new infrastructure development and expansion in developing nations, we study infrastructure maintenance in a developed economy. Focusing on existing roads with pre-determined renewal period and narrowing down on close elections helps us avoid endogeneity issues coming from placement and the timing of new infrastructure. Our approach provides insights that are more relevant for policymakers in advanced economies, where road infrastructure is well established but requires continuous upkeep. 
Second, we compute a novel measure of road quality using a fine-tuned vision model trained on satellite images. Current literature has relied on measures based on Pavement Condition Index (PCI) or International Roughness Index (IRI) to measure road quality, but these are not available for many localities. Similarly, \cite{currier2023} measures road roughness using vertical acceleration data from a smartphone app, but this also requires large-scale data available mainly for large metropolitan areas and is not available for smaller towns. Our approach allows us to measure road quality at the local level, bypassing the limitations of existing measures. 
Third, we focus on long-term effects of reduced road taxes funded mainly via property taxes, and observe statistically significant decline in house prices occurring after four years and persisting in later years. This delayed effect highlights how the consequences of road tax cuts are not immediate but gradually accumulate as funding is reduced and roads deteriorate. Our research also identifies heterogeneous impacts, showing that urban areas and higher-priced houses suffer more pronounced price declines. By focusing on distributional effects, the study adds nuance to the understanding of how infrastructure maintenance affects local housing markets.

% Third, while much of the literature focuses on new infrastructure development and expansion in developing nations, we study infrastructure maintenance in a developed economy.  Focusing on existing roads mitigates an important source of identification problems. Our approach provides insights that are more relevant for policymakers in advanced economies, where road infrastructure is well established but requires continuous upkeep. 

% These results are consistent with other work on impact of roads on house prices. For instance, \cite{gonzalez2016paving} find that paving roads in Acayucan, Mexico, increases property values by 28\%, while \cite{theisen2021road} find effect of 13\% in nearest towns where roads were paved in Norway.

